%% ============================================================
%%  Resumen y Abstract
%%  Duenio inicial: claude-1
%%  Reglas: ver ../../CLAUDE.md
%%  No editar si TASKS.md marca este archivo IN_PROGRESS por otro agente.
%% ============================================================

\chapter*{Resumen}
\addcontentsline{toc}{chapter}{Resumen/Abstract\vspace{-0.5cm}}

Los modelos profundos basados en atención superan a los modelos clásicos en desempeño predictivo sobre problemas multivariables, pero no entregan una descripción de las relaciones entre variables que sostiene esa predicción. Las técnicas de explicabilidad establecidas tampoco la entregan, porque asignan importancia a cada variable por separado o dibujan la relevancia como un mapa continuo, y ninguno de esos objetos puede enunciarse fuera del modelo que lo produjo, compararse con conocimiento previo del fenómeno ni someterse a refutación. De ahí surge la pregunta que ordena el trabajo: si es posible transformar las matrices de atención de un modelo Transformer en grafos dirigidos compactos que generen hipótesis relacionales reproducibles y falsables, con mayor fidelidad al comportamiento del modelo que las técnicas tradicionales.

Para responderla se propone un framework metodológico que opera sobre un modelo ya entrenado, sin reentrenarlo ni modificar su arquitectura, y que encadena tres transformaciones: convertir la matriz de atención en un grafo dirigido sobre variables con nombre, reducir ese grafo a una hipótesis relacional compacta y evaluar esa hipótesis contra el comportamiento del modelo. La contribución no está en extraer el grafo, operación algorítmicamente sencilla que varios trabajos ya resuelven, sino en el criterio que decide cuándo una relación extraída merece tratarse como hipótesis científica. Ese criterio se organiza como una cadena de evidencia acumulativa, donde cada relación debe mostrar que reaparece al reentrenar el modelo, que resiste el cambio de la regla con que se discretiza, que es improbable bajo una hipótesis nula explícita, que el modelo efectivamente depende de ella, y que su efecto es específico. Los umbrales quedan declarados antes de ejecutar el procedimiento, y la evidencia se contrasta con paradigmas que no comparten supuestos.

Este informe corresponde a la etapa de avance y documenta el diseño, no su ejecución. El framework está formalizado, el criterio de validación está especificado y el estado del arte que lo sustenta proviene de una revisión sistemática conducida bajo el protocolo PRISMA 2020, de la que salieron las 25 referencias que respaldan el trabajo. Aplicar el procedimiento sobre el caso de estudio, un conjunto de índices espectrales derivados de imágenes satelitales, es el programa de la etapa siguiente y aparece calendarizado en el plan de trabajo.

{\bf Palabras clave:} interpretabilidad, mecanismos de atención, grafo relacional, series de tiempo multivariadas, revisión sistemática.

\vspace{2cm}
\begingroup
\let\clearpage\relax
\chapter*{Abstract}

Attention-based deep models outperform classical models on multivariate prediction tasks, yet they do not expose the relations among variables that support those predictions. Established explainability techniques do not expose them either, since they assign importance to each variable in isolation or render relevance as a continuous map, and neither object can be stated outside the model that produced it, compared against prior knowledge of the phenomenon, or subjected to refutation. Hence the question that organises this work: whether the attention matrices of a Transformer model can be turned into compact directed graphs that yield reproducible and falsifiable relational hypotheses, with greater faithfulness to the model's behaviour than traditional techniques offer.

The proposed answer is a methodological framework that operates on an already trained model, without retraining it or modifying its architecture, and that chains three transformations: turning the attention matrix into a directed graph over named variables, reducing that graph to a compact relational hypothesis, and evaluating that hypothesis against the behaviour of the model. The contribution does not lie in extracting the graph, which is algorithmically simple and already solved in several works, but in the criterion that decides when an extracted relation deserves to be treated as a scientific hypothesis. That criterion is organised as a cumulative chain of evidence, in which each relation must show that it reappears when the model is retrained, that it withstands a change in the discretisation rule, that it is unlikely under an explicit null hypothesis, that the model actually depends on it, and that its effect is specific. Thresholds are declared before the procedure runs, and the evidence is contrasted with paradigms that do not share assumptions.

This document reports on the progress stage and describes the design rather than its execution. The framework is formalised, the validation criterion is specified, and the supporting state of the art comes from a systematic review conducted under the PRISMA 2020 protocol, which yielded the 25 references behind this work. Running the procedure on the case study, a set of spectral indices derived from satellite imagery, is the programme for the next stage and appears scheduled in the work plan.

{\bf Keywords:} interpretability, attention mechanisms, relational graph, multivariate time series, systematic review.
\endgroup
